% =====================================================================
% Understanding RWA, Impairment, and Risk-Sensitive Pricing
% A Study Textbook on the A-IRB Risk-Weight Function, IFRS 9 Expected
% Credit Losses, and Risk-Based Loan Pricing
%
% Compile with:  lualatex Understanding_RWA_Impairment_Pricing.tex
%   (run twice for the table of contents and cross-references;
%    pdflatex also works, lualatex is recommended for the fonts)
%
% Companion scripts: RWA_visualised.py, Impairment_visualised.py,
%                    Pricing_visualised.py
%
% Created with AI assistance; see Chapter 1 for provenance and
% disclaimers.
% =====================================================================
\documentclass[11pt,a4paper]{report}

% ---------- geometry -------------------------------------------------
\usepackage[a4paper,margin=2.5cm,headheight=14pt]{geometry}

% ---------- typography (matching the author's TeX setup) --------------
\usepackage{tgpagella}          % Palatino-like text font (as used by moderncv)
\usepackage{newpxmath}          % matching math font
\usepackage[T1]{fontenc}
\usepackage{microtype}

% ---------- math ------------------------------------------------------
\usepackage{amsmath}
\usepackage{amssymb}
\usepackage{amsthm}
\usepackage{mathtools}

% ---------- tables, figures ------------------------------------------
\usepackage{booktabs}
\usepackage{array}
\usepackage{tikz}
\usetikzlibrary{arrows.meta,positioning}

% ---------- colours (echoing the cerulean scheme of the author's CV) ---
\usepackage{xcolor}
\usepackage{eurosym}
\definecolor{ceruleanlink}{HTML}{007B8A}
\definecolor{accentgrey}{HTML}{555555}
\newcommand{\EUR}[1]{\euro#1}

% ---------- links ----------------------------------------------------
\usepackage[colorlinks=true,linkcolor=ceruleanlink,citecolor=ceruleanlink,urlcolor=ceruleanlink]{hyperref}
\usepackage{enumitem}

% ---------- headers ----------------------------------------------------
\usepackage{fancyhdr}
\pagestyle{fancy}
\fancyhf{}
\fancyhead[LE]{\small\itshape Understanding RWA, Impairment, and Pricing}
\fancyhead[RO]{\small\itshape\nouppercase{\leftmark}}
\fancyfoot[C]{\thepage}
\renewcommand{\headrulewidth}{0.4pt}

% ---------- theorem environments ---------------------------------------
\theoremstyle{definition}
\newtheorem{definition}{Definition}[chapter]
\newtheorem{example}[definition]{Example}
\newtheorem{remark}[definition]{Remark}
\theoremstyle{plain}
\newtheorem{lemma}[definition]{Lemma}
\newtheorem{theorem}[definition]{Theorem}
\newtheorem{corollary}[definition]{Corollary}

% ---------- convenience macros ------------------------------------------
\newcommand{\PD}{\mathrm{PD}}
\newcommand{\LGD}{\mathrm{LGD}}
\newcommand{\EAD}{\mathrm{EAD}}
\newcommand{\RWA}{\mathrm{RWA}}
\newcommand{\EL}{\mathrm{EL}}
\newcommand{\EIR}{\mathrm{EIR}}
\newcommand{\ECL}{\mathrm{ECL}}
\newcommand{\raroc}{\mathrm{RAROC}}
\newcommand{\keyinsight}[1]{\par\vspace{2mm}\noindent\fbox{\parbox{0.96\textwidth}{\textbf{Key insight.} #1}}\par\vspace{2mm}}

\title{\textbf{Understanding RWA, Impairment, and\\[2mm] Risk-Sensitive Pricing}\\[6mm]
\large A Study Textbook on the A-IRB Risk-Weight Function,\\ IFRS~9 Expected Credit Losses, and Risk-Based Loan Pricing}
\author{}
\date{October 2026}

\begin{document}

\maketitle
\thispagestyle{empty}
\vfill
\begin{center}
\itshape Running example throughout this book:\\
an unsecured consumer instalment loan of \EUR{10{,}000}.
\end{center}
\vfill

\tableofcontents

% =====================================================================
\chapter{Introduction, Provenance, and How to Use This Book}
% =====================================================================

\section{Purpose and scope}

A bank that grants a loan answers three quantitatively different
questions about the same credit risk:

\begin{enumerate}[label=(\alph*)]
  \item \textbf{Regulation (CRR/Basel):} \emph{How much own funds must we
        hold against the \emph{unexpected} loss of this loan?} The answer
        is expressed in risk-weighted assets (RWA) and derived from the
        A-IRB risk-weight function --- the \emph{Gordy formula}.
  \item \textbf{Accounting (IFRS 9):} \emph{How much \emph{expected} loss
        must we recognise as an impairment allowance today?} The answer is
        the expected credit loss (ECL), possibly probability-weighted
        across macroeconomic scenarios.
  \item \textbf{Management (pricing):} \emph{What interest rate must we
        charge so that the loan earns its cost of capital after covering
        funding, operating cost, and expected loss?} The answer is the
        risk-sensitive break-even rate derived from a RAROC condition.
\end{enumerate}

This book develops all three answers from first principles, at the level
of a business-mathematics graduate (probability, statistics, financial
mathematics): short proofs where feasible, proof sketches where a full
proof would exceed the scope, worked numerical examples in every chapter,
and references to the primary regulatory and academic sources. One single
loan --- an unsecured consumer instalment loan of \EUR{10{,}000} ---
serves as the running example throughout, so that the three frameworks
can be compared on identical numbers.

Three companion Python scripts (\texttt{RWA\_visualised.py},
\texttt{Impairment\_visualised.py}, \texttt{Pricing\_visualised.py})
implement the mathematics of Parts~A--C as interactive Flask dashboards;
Chapter~\ref{chap:dashboards} explains how to run them and how every
number in this book can be reproduced with them.

\section{How this book was created}

This textbook was drafted by a large language model (AI) at the request
of the reader, in an iterative process: the mathematical content follows
published regulatory and academic sources (cited in the References and
rated in the Source Notes); every numerical example was computed with
executable code, not asserted; and two errors found in an earlier draft
(an incorrect risk weight for the PD~$=8\%$ comparison and an
inconsistent lifetime-EL pricing basis) were corrected against
computation before publication. The prose was then revised for
textbook-style exposition.

\textbf{Disclaimers.} The reader should treat this book accordingly:

\begin{enumerate}
  \item This is an educational treatment with deliberate didactic
        simplifications. It is not legal, regulatory, accounting, tax,
        or investment advice.
  \item The example parameters (PD~$=2\%$, LGD~$=45\%$, \dots) are
        illustrative. A real A-IRB bank's parameters come from internally
        validated, supervisor-approved models and cannot be reconstructed
        from outside the bank.
  \item The pricing model is a teaching model. Real banks add acquisition
        costs, liquidity costs, Pillar~2 add-ons, MREL and
        leverage-ratio constraints, competitive considerations, and
        behavioural effects.
  \item Regulation changes. CRR amendments (``CRR3'', output floors) and
        IFRS~9 interpretation practice evolve; always check the current
        consolidated text on EUR-Lex and the IFRS Foundation website.
  \item Nothing in this book predicts whether any specific borrower
        receives a loan, or at what rate. The book explains how a bank
        \emph{measures and prices} credit risk, not whether credit
        \emph{is granted}.
\end{enumerate}

\section{Notation}

$\Phi$ denotes the standard normal distribution function and
$G = \Phi^{-1}$ its inverse (quantile function). $\mathbb{E}$ is
expectation, $\mathbb{P}$ probability, $\Var$ variance. Currency amounts
are euros. A full notation table is given in
Appendix~\ref{app:notation}.

% =====================================================================
\chapter{Mathematical Foundations}
\label{chap:foundations}
% =====================================================================

This chapter collects the probabilistic and financial tools used later,
so that the derivations in Parts~A--C stay self-contained. Readers with
a statistics background may skim it and return on demand.

\section{The standard normal distribution and its quantiles}

\begin{definition}[Standard normal distribution]
A random variable $Z$ is \emph{standard normally distributed}
($Z\sim\mathcal{N}(0,1)$) if its density is
\begin{equation}
  \varphi(z)=\frac{1}{\sqrt{2\pi}}\,e^{-z^{2}/2},\qquad z\in\mathbb{R},
\end{equation}
with distribution function
\begin{equation}
  \Phi(x)=\int_{-\infty}^{x}\varphi(z)\,dz .
\end{equation}
\end{definition}

Two numerical values recur throughout this book and should be memorised:
\begin{equation}
  G(0.02) \approx -2.0537,
  \qquad
  G(0.999) \approx 3.0902 .
\end{equation}
The first says: a $2\%$ default probability corresponds to a default
threshold $2.05$ standard deviations \emph{below} the mean of the
credit-worthiness variable. The second says: the 99.9\%\ quantile sits
$3.09$ standard deviations \emph{above} the mean of the systematic
factor --- the supervisory stress state.

\begin{lemma}[Linear combinations of independent normals]
\label{lem:linnorm}
Let $Y,\varepsilon$ be independent with
$Y,\varepsilon\sim\mathcal{N}(0,1)$ and let $a,b\in\mathbb{R}$. Then
\begin{equation}
  X=aY+b\varepsilon \sim \mathcal{N}\!\bigl(0,\; a^{2}+b^{2}\bigr).
\end{equation}
\end{lemma}

\begin{proof}
The characteristic function of $X$ is
\begin{equation}
  \mathbb{E}\bigl[e^{it(aY+b\varepsilon)}\bigr]
  =\mathbb{E}\bigl[e^{itaY}\bigr]\,\mathbb{E}\bigl[e^{itb\varepsilon}\bigr]
  =e^{-(at)^{2}/2}\,e^{-(bt)^{2}/2}
  =e^{-(a^{2}+b^{2})t^{2}/2},
\end{equation}
using independence and the characteristic function
$e^{-t^{2}/2}$ of the standard normal. The result is the characteristic
function of a $\mathcal{N}(0,a^{2}+b^{2})$ variable, and
characteristic functions determine the distribution uniquely.
\end{proof}

\section{Conditional probability and conditional expectation}

For events $A$ and $B$ with $\mathbb{P}(B)>0$,
\begin{equation}
  \mathbb{P}(A\mid B)=\frac{\mathbb{P}(A\cap B)}{\mathbb{P}(B)} .
\end{equation}
If $X=aY+b\varepsilon$ and we \emph{condition on} the realised value
$Y=y$, then $X$ becomes a deterministic shift plus noise:
\begin{equation}
  X \mid Y=y \;\sim\; \mathcal{N}\!\bigl(ay,\; b^{2}\bigr).
  \label{eq:condnormal}
\end{equation}
Equation~\eqref{eq:condnormal} is the engine of the entire A-IRB
derivation: conditioning on a bad economy shifts the mean of every
borrower's credit-worthiness downward and thereby raises every
borrower's default probability.

\section{Expectation, variance decomposition, and tail risk}

For a loss variable $L\ge 0$,
\begin{equation}
  \EL=\mathbb{E}[L] \qquad\text{(expected loss)},
\end{equation}
while economic or regulatory capital addresses the \emph{upper tail}:
\begin{equation}
  \text{capital} \approx \text{tail quantile of } L \;-\; \EL .
\end{equation}
This ``tail minus expectation'' logic is the single most important
structural idea of this book:

\keyinsight{Expected loss is a \emph{cost of doing business} --- it is
priced into the interest rate and provisioned under IFRS~9.
\emph{Unexpected} loss --- the deviation beyond expectation in a severe
scenario --- is what regulatory capital must absorb. Part~A computes
the tail term, Part~B the expectation term, and Part~C charges both to
the customer.}

\section{Present value and the annuity}

A cash flow $C$ received at time $t$ has present value
$C\,(1+r)^{-t}$ at discount rate $r$. For a loan of principal $N$,
nominal rate $r$, and term $n$ years with constant annual annuity $A$,
the condition $\PV= N$ gives
\begin{equation}
  A=\frac{N\,r}{1-(1+r)^{-n}} .
  \label{eq:annuity}
\end{equation}
\begin{example}[Running example, annuity]
With $N=\EUR{10{,}000}$, $r=8.5\%$, $n=3$:
\begin{equation}
  A=\frac{10{,}000\times 0.085}{1-1.085^{-3}}=\EUR{3{,}915.39} .
\end{equation}
The amortisation schedule (used in Parts~B and~C) is shown in
Table~\ref{tab:amortisation}.
\end{example}

\begin{table}[htbp]
\centering
\caption{Amortisation schedule of the running example
(\EUR{10{,}000}, $8.5\%$, 3 years, annuity \EUR{3{,}915.39}).}
\label{tab:amortisation}
\begin{tabular}{rrrrr}
\toprule
Year & Balance (BoY) & Interest & Principal & Balance (EoY)\\
\midrule
1 & 10{,}000.00 & 850.00 & 3{,}065.39 & 6{,}934.61\\
2 & 6{,}934.61 & 589.44 & 3{,}325.95 & 3{,}608.66\\
3 & 3{,}608.66 & 306.74 & 3{,}608.66 & 0.00\\
\bottomrule
\end{tabular}
\end{table}

% =====================================================================
\chapter{Part A --- Risk-Weighted Assets under the A-IRB Approach}
\label{chap:rwa}
% =====================================================================

\section{The regulatory question}

\begin{figure}[htbp]
\centering
\begin{tikzpicture}[node distance=10mm,
  box/.style={draw=ceruleanlink, rounded corners, fill=white,
              minimum width=52mm, minimum height=8mm, align=center}]
\node[box] (params) {Risk parameters $\PD$, $\LGD$, $\EAD$};
\node[box, below=of params] (rwf) {Risk-weight function (Gordy): capital factor $K$};
\node[box, below=of rwf] (rwa) {$\RWA = 12.5 \times K \times \EAD$};
\node[box, below=of rwa] (cap) {Own funds requirement $=8\%\times\RWA$};
\draw[-{Latex[length=2mm]}] (params) -- (rwf);
\draw[-{Latex[length=2mm]}] (rwf) -- (rwa);
\draw[-{Latex[length=2mm]}] (rwa) -- (cap);
\end{tikzpicture}
\caption{The A-IRB computation chain for a single exposure.}
\label{fig:chain}
\end{figure}

Pillar~1 of the capital framework requires banks to hold own funds of at
least $8\%$ of their risk-weighted assets. Two building blocks therefore
have to be understood: \emph{risk-weighted assets} and the $8\%$\ rule.

\begin{definition}[Risk-weighted assets]
The risk-weighted assets of an exposure are its exposure value
$\EAD$ multiplied by a \emph{risk weight} $RW$:
\begin{equation}
  \RWA = RW \times \EAD .
\end{equation}
In the standardised approach $RW$ is a fixed supervisory percentage;
in the IRB approaches $RW$ is \emph{computed} from risk parameters.
\end{definition}

\section{The factor 12.5: a complete derivation}

\begin{theorem}[RWA scaling]
Let the minimum own funds ratio be $8\%$, and let the capital
requirement for one exposure be $K\times\EAD$ (with $K$ derived from the
risk-weight function). Then the risk-weighted assets that make the two
statements consistent are
\begin{equation}
  \RWA = \frac{K\times\EAD}{0.08} = 12.5 \times K \times \EAD .
  \label{eq:scaling}
\end{equation}
\end{theorem}

\begin{proof}
By definition of the capital ratio, the requirement is
\begin{equation}
  \text{own funds} = 0.08 \times \RWA .
\end{equation}
The IRB capital requirement for the exposure is
$\text{own funds} = K \times \EAD$. Equating and dividing by $0.08$
($1/0.08 = 12.5$) yields \eqref{eq:scaling}.
\end{proof}

This is pure algebra --- there is no economics in the $12.5$; the economics
sits entirely inside $K$, to which we now turn.

\section{The three risk parameters}

\begin{description}
\item[PD (probability of default).] The probability that the borrower
      defaults within one year under the supervisory default definition
      (CRR Art.~178: 90 days past due or unlikeliness to pay). Running
      example: $\PD=2\%$.
\item[LGD (loss given default).] The expected loss as a fraction of the
      exposure \emph{given} default, net of recoveries, including
      workout costs and discounting of recoveries (CRR Art.~181).
      Unsecured consumer loans have no collateral, but collections
      recoveries still reduce the loss. Running example: $\LGD=45\%$.
\item[EAD (exposure at default).] The expected outstanding amount at the
      moment of default. For a fully drawn instalment loan the EAD
      approximates the outstanding balance; immediately after
      disbursement, $\EAD\approx\EUR{10{,}000}$. Under A-IRB the bank
      estimates EAD internally.
\end{description}

\begin{remark}[F-IRB versus A-IRB]
Under the Foundation IRB approach the bank estimates only PD; LGD and
EAD follow supervisory values. Under the \emph{Advanced} IRB approach
the bank estimates PD, LGD, and EAD with internal, supervisor-approved
models. The risk-weight function itself is prescribed --- the ``A''
refers to the parameter estimation, not to the formula.
\end{remark}

\section{The single-factor model: from Merton to Vasicek}

The IRB formula is not an arbitrary regulatory curve. Its intellectual
pedigree is Merton (1974) $\to$ Vasicek (2002) $\to$ Gordy (2003)
$\to$ Basel II/III IRB $\to$ CRR Arts.~153--154. We derive it in four
steps.

\subsection{Step 1: the factor model}

\begin{definition}[Asymptotic single risk factor model]
For each borrower $i$, the standardised asset (credit-worthiness) return
is
\begin{equation}
  X_i=\sqrt{R}\;Y+\sqrt{1-R}\;\varepsilon_i ,
  \label{eq:asrf}
\end{equation}
where $Y\sim\mathcal{N}(0,1)$ is a single systematic factor (the
economy), $\varepsilon_i\sim\mathcal{N}(0,1)$ is idiosyncratic noise,
all variables mutually independent, and $R\in(0,1)$ is the asset
correlation.
\end{definition}

\begin{lemma}[$X_i$ is standard normal]
\label{lem:standardnormal}
$X_i\sim\mathcal{N}(0,1)$.
\end{lemma}

\begin{proof}
By Lemma~\ref{lem:linnorm} with $a=\sqrt R$, $b=\sqrt{1-R}$:
$\Var(X_i)=R+(1-R)=1$ and mean $0$; normality follows from the linear
combination of independent Gaussians.
\end{proof}

Interpretation: a borrower's fortunes are driven by the macroeconomy
(weight $\sqrt R$) and by personal circumstances (weight
$\sqrt{1-R}$). The correlation of any two borrowers' returns is
\begin{equation}
  \operatorname{Cov}(X_i,X_j)=R\Var(Y)=R,
\end{equation}
which explains the name ``asset correlation'' for $R$.

\subsection{Step 2: the default threshold}

\begin{definition}[Merton default rule]
Borrower $i$ defaults in the period iff $X_i<c$, where the threshold $c$
is chosen such that the \emph{unconditional} default probability equals
the PD:
\begin{equation}
  \mathbb{P}(X_i<c)=\PD .
\end{equation}
\end{definition}

Since $X_i\sim\mathcal{N}(0,1)$, the threshold is
\begin{equation}
  \Phi(c)=\PD
  \qquad\Longleftrightarrow\qquad
  c=G(\PD)=\Phi^{-1}(\PD).
  \label{eq:threshold}
\end{equation}
For $\PD=2\%$: $c=G(0.02)=-2.0537$. Default occurs $2.05$ standard
deviations below average fortune --- a rare event in good times.

\subsection{Step 3: the conditional default probability}

Now condition on the economy being in a specific state $Y=y$. Using
\eqref{eq:condnormal} with $a=\sqrt R$, $b=\sqrt{1-R}$:
\begin{equation}
  X_i\mid Y=y \;\sim\; \mathcal{N}\!\bigl(\sqrt{R}\,y,\;1-R\bigr),
\end{equation}
hence
\begin{align}
  \mathbb{P}(\text{default}\mid Y=y)
  &=\mathbb{P}\!\bigl(\sqrt{R}\,y+\sqrt{1-R}\,\varepsilon<G(\PD)\bigr)\\
  &=\mathbb{P}\!\left(\varepsilon<
      \frac{G(\PD)-\sqrt{R}\,y}{\sqrt{1-R}}\right)
   =\Phi\!\left(\frac{G(\PD)-\sqrt{R}\,y}{\sqrt{1-R}}\right).
  \label{eq:condpd}
\end{align}

The regulator fixes the adverse systematic state at the 99.9\%\ quantile.
With the CRR sign convention (adverse state in the upper tail),
\begin{equation}
  y^{\text{stress}}=G(0.999)\approx 3.0902 ,
\end{equation}
and the \emph{stress default probability} becomes
\begin{equation}
  \boxed{\;
  \PD^{\text{stress}}
  =\Phi\!\left[
  \frac{1}{\sqrt{1-R}}\,G(\PD)
  +\sqrt{\frac{R}{1-R}}\;G(0.999)
  \right]}
  \label{eq:pdstress}
\end{equation}
--- the term in square brackets is simply
\eqref{eq:condpd} with $y=G(0.999)$, split over the common denominator
$\sqrt{1-R}$.

\begin{remark}[Why the stress PD can be many times the PD]
The correlation term multiplies $G(0.999)\approx 3.09$ by
$\sqrt{R/(1-R)}$. For $R\approx 9.5\%$ this contributes roughly $+1.0$
to the argument of $\Phi$. Near the far left tail of $\Phi$, one unit
of argument is worth a \emph{multiplication} of the tail probability:
with $\PD=2\%$ the stress PD is $12.3\%$ --- six times higher. This
systematic amplification is precisely the economic content of the
correlation.
\end{remark}

\subsection{Step 4: capital equals stress loss minus expected loss}

The loss rate of the exposure, \emph{given} default, is $\LGD$. In the
stress state the expected loss rate is therefore
$\LGD\times\PD^{\text{stress}}$; the unconditional expected loss rate is
$\EL=\PD\times\LGD$. Capital covers the \emph{unexpected} part:

\begin{equation}
  \boxed{\;
  K=\LGD\times\Phi\!\left[
  \frac{1}{\sqrt{1-R}}\,G(\PD)
  +\sqrt{\frac{R}{1-R}}\;G(0.999)\right]-\PD\times\LGD}
  \label{eq:gordy}
\end{equation}

This is the \textbf{Gordy formula}. Factoring out $\LGD$ gives the
transparent form $K=\LGD\times(\PD^{\text{stress}}-\PD)$: capital
equals the loss-given-default times the \emph{excess} of the stress
default probability over the ordinary one.

\section{Why the formula applies loan-by-loan: Gordy's theorem}

Regulation applies \eqref{eq:gordy} exposure by exposure. That is only
legitimate because of a remarkable asymptotic property.

\begin{theorem}[Gordy (2003), portfolio invariance --- proof sketch]
\label{thm:gordy}
Consider a portfolio of $n$ exposures of equal size, each following
\eqref{eq:asrf} with the same $R$ and independent
$\varepsilon_i$. Conditionally on $Y$, defaults are independent, each
with probability $\PD^{\text{stress}}(Y)$. By the law of large numbers,
as $n\to\infty$ the portfolio default rate converges to
$\PD^{\text{stress}}(Y)$ almost surely; idiosyncratic risk vanishes. The
portfolio loss rate at the 99.9\%\ quantile of $Y$ is therefore
exactly $\LGD\times\PD^{\text{stress}}$, with no residual idiosyncratic
term, and the capital per unit of exposure depends \emph{only} on the
exposure's own $(\PD,\LGD)$ and on the common $R$ --- not on the
portfolio's composition.
\end{theorem}

The full proof (Gordy 2003) requires the uniform-exposure-size and
single-factor assumptions; both are idealisations. Real portfolios are
finite and concentrated, which is why banks hold granularity buffers and
why economic capital (Chapter~\ref{chap:pricing}) exceeds regulatory
capital. But the theorem explains the structural feature that makes the
IRB usable: \emph{capital is additive because idiosyncratic risk is
diversified away in the limit, and what remains is systematic risk,
priced identically for every exposure with the same parameters.}

\section{The CRR retail correlation function}

For retail exposures, the CRR (Art.~154) does not let the bank choose
$R$ freely. It prescribes a decreasing function of the PD. Writing
\begin{equation}
  A(\PD)=\frac{1-e^{-35\,\PD}}{1-e^{-35}},
\end{equation}
the correlation is
\begin{equation}
  \boxed{\;
  R(\PD)=0.03\times A(\PD)+0.16\times\bigl(1-A(\PD)\bigr)
        =0.16-0.13\times A(\PD)}
  \label{eq:corr}
\end{equation}

Interpretation of the two anchors:

\begin{itemize}
\item $\PD\to 0$: $A\to 0$, hence $R\to 16\%$. Very safe borrowers are
      driven almost entirely by long macro cycles --- think of prime
      mortgage borrowers who only default in system-wide crises.
\item $\PD\to 1$: $A\to 1$, hence $R\to 3\%$. Already-distressed
      borrowers default for idiosyncratic reasons; the economy matters
      less for \emph{whether} they default.
\end{itemize}

For special retail sub-classes the CRR replaces \eqref{eq:corr} by
constants: $R=0.15$ for exposures secured by immovable property, and
$R=0.04$ for qualifying revolving retail exposures. Our unsecured
instalment loan uses the general function \eqref{eq:corr}.

\section{Worked example: the \EUR{10{,}000} loan, step by step}
\label{sec:workedA}

Assumptions (illustrative): $\EAD=\EUR{10{,}000}$, $\PD=2\%$,
$\LGD=45\%$, non-revolving unsecured retail.

\paragraph{Step 1 --- Correlation.}
With $35\times\PD = 0.7$:
\begin{align}
  A &= \frac{1-e^{-0.7}}{1-e^{-35}}=0.5034,\\
  R &= 0.16-0.13\times 0.5034 = 0.0946 \quad(9.46\%).
\end{align}

\paragraph{Step 2 --- Quantiles.}
\begin{equation}
  G(0.02)=-2.0537,\qquad G(0.999)=3.0902 .
\end{equation}

\paragraph{Step 3 --- Stress argument.}
\begin{equation}
  z=\frac{-2.0537}{\sqrt{1-0.0946}}
     +\sqrt{\frac{0.0946}{1-0.0946}}\times 3.0902
    =-2.1561+1.0000=-1.1597 .
\end{equation}

\paragraph{Step 4 --- Stress PD.}
\begin{equation}
  \PD^{\text{stress}}=\Phi(-1.1597)=12.31\% .
\end{equation}
A once-in-a-thousand-years downturn raises this borrower's default
probability from $2\%$ to $12.3\%$.

\paragraph{Step 5 --- Capital factor.}
\begin{equation}
  K=0.45\times(0.1231-0.02)=0.45\times 0.1031=0.0464\quad(4.64\%).
\end{equation}

\paragraph{Step 6 --- RWA and own funds.}
\begin{align}
  RW &= 12.5\times K = 58.0\%,\\
  \RWA &= 0.5799\times 10{,}000 = \EUR{5{,}798.64},\\
  \text{own funds} &= 0.08\times 5{,}798.64=\EUR{463.89}.
\end{align}

\begin{table}[htbp]
\centering
\caption{Part A summary for the running example.}
\begin{tabular}{lr}
\toprule
Quantity & Value\\
\midrule
$\PD$ / $\LGD$ / $\EAD$ & $2\%$ / $45\%$ / \EUR{10{,}000}\\
Correlation $R$ & $9.46\%$\\
Stress PD (99.9\%) & $12.31\%$\\
Capital factor $K$ & $4.64\%$\\
Risk weight $RW$ & $58.0\%$\\
$\RWA$ & \EUR{5{,}798.64}\\
Expected loss $\PD\times\LGD\times\EAD$ & \EUR{90.00}\\
Own funds requirement ($8\%\times\RWA$) & \EUR{463.89}\\
\bottomrule
\end{tabular}
\end{table}

\keyinsight{The bank neither expects to lose \EUR{5{,}798.64} nor
\EUR{463.89}. It \emph{expects} \EUR{90} of loss (covered by pricing and
provisions) and must \emph{hold} \EUR{463.89} of capital against the
99.9\%\ tail deviation beyond that expectation.}

\section{Sensitivities}

Table~\ref{tab:pdsens} varies the PD; Table~\ref{tab:lgdsens} the LGD.
The required rate column anticipates Chapter~\ref{chap:pricing}
($12$-month EL basis, funding $2\%$, opex \EUR{166.67} p.a., hurdle
$12\%$, regulatory capital).

\begin{table}[htbp]
\centering
\caption{PD sensitivity ($\LGD=45\%$, $\EAD=\EUR{10{,}000}$).}
\label{tab:pdsens}
\begin{tabular}{rrrrrrrr}
\toprule
$\PD$ & $R$ & $\PD^{\text{stress}}$ & $K$ & $RW$ & $\RWA$ & own funds & req.\ rate\\
\midrule
0.5\% & 13.9\% & 6.25\% & 2.59\% & 32.4\% & 3{,}236 & 258.89 & 5.22\%\\
1.0\% & 12.2\% & 9.14\% & 3.66\% & 45.8\% & 4{,}577 & 366.18 & 5.73\%\\
2.0\% & 9.5\% & 12.31\% & 4.64\% & 58.0\% & 5{,}799 & 463.89 & 6.56\%\\
5.0\% & 5.3\% & 16.81\% & 5.31\% & 66.4\% & 6{,}642 & 531.32 & 8.65\%\\
8.0\% & 3.8\% & 20.64\% & 5.69\% & 71.1\% & 7{,}108 & 568.64 & 10.69\%\\
15\% & 3.1\% & 30.75\% & 7.09\% & 88.6\% & 8{,}860 & 708.81 & 15.53\%\\
30\% & 3.0\% & 50.44\% & 9.20\% & 115.0\% & 11{,}498 & 919.82 & 25.76\%\\
\bottomrule
\end{tabular}
\end{table}

\begin{table}[htbp]
\centering
\caption{LGD sensitivity ($\PD=2\%$, $\EAD=\EUR{10{,}000}$).}
\label{tab:lgdsens}
\begin{tabular}{rrrr}
\toprule
$\LGD$ & $K$ & $\RWA$ & $\EL$ (p.a.)\\
\midrule
10\% & 1.03\% & 1{,}288.59 & 20.00\\
25\% & 2.58\% & 3{,}221.47 & 50.00\\
45\% & 4.64\% & 5{,}798.64 & 90.00\\
65\% & 6.70\% & 8{,}375.82 & 130.00\\
90\% & 9.28\% & 11{,}597.29 & 180.00\\
\bottomrule
\end{tabular}
\end{table}

\begin{remark}[Reading the PD table]
$K$ rises with $\PD$ but \emph{less than proportionally}, because the
correlation falls as $\PD$ rises ($16\%\to 3\%$), cushioning the stress
amplification. The risk weight is thus concave in $\PD$. RWA is
linear in $\EAD$ and $\LGD$. The pricing-relevant quantity, however, is
dominated by expected loss: between $\PD=2\%$ and $\PD=8\%$ the own
funds requirement grows by only $23\%$ while the EL grows by a factor
of four --- a first glimpse of why risk-based pricing is driven by
Part~B more than by Part~A.
\end{remark}

\section{Limitations of Part A}

\begin{itemize}
\item Real A-IRB banks face PD floors, downturn-LGD requirements,
      CCF/EAD model risk, and (since CRR2) an output floor tying IRB
      RWA to a percentage of standardised RWA --- the applicable
      percentage depends on the bank's transition schedule and current
      CRR state.
\item The ASRF assumptions (infinite granularity, single factor) fail
      for concentrated portfolios; supervisory add-ons and Pillar~2
      address this.
\item The 99.9\%\ level is a calibration choice, not an empirical fact
      about economies.
\end{itemize}

% =====================================================================
\chapter{Part B --- Impairment under IFRS 9}
\label{chap:ifrs9}
% =====================================================================

\section{The accounting question}

Accounting asks a different question than regulation. Regulation asks
how much capital to hold against unexpected loss; IFRS~9 asks how much
\emph{expected} loss to recognise \emph{today} as an allowance
(impairment) in profit or loss. IFRS~9 (effective 2018) replaced the
IAS~39 incurred-loss model --- under which losses were recognised only
when a triggering event had occurred --- with a forward-looking
expected-credit-loss model. The philosophical change: loss recognition
moved from \emph{evidence of damage} to \emph{probability-weighted
forecast}.

\section{The three-stage model}

IFRS~9 (paras 5.5.15--5.5.20) sorts financial assets into three stages.
The classification governs two things: the \emph{measurement horizon}
of the allowance and the \emph{interest recognition basis}.

\begin{center}
\begin{tabular}{lll}
\toprule
Stage & Trigger & Allowance / interest basis\\
\midrule
1 & no significant deterioration & 12-month ECL; interest gross\\
2 & significant increase in credit risk & lifetime ECL; interest gross\\
3 & credit-impaired (default-like) & lifetime ECL; interest net\\
\bottomrule
\end{tabular}
\end{center}

The 12-month/lifetime distinction is the \emph{cliff} of the model.
Crossing a stage-2 trigger multiplies the allowance for a multi-period
loan, which is why stage migration --- not default itself --- is a major
driver of bank P\&L volatility in downturns.

\begin{definition}[Significant increase in credit risk (SICR)]
A stage-2 migration occurs when the lifetime PD at reporting rises
significantly above the lifetime PD expected at initial recognition.
Common operationalisations: 30+ days past due (rebuttable
presumption), internal rating deterioration by a defined number of
notches, or lifetime PD more than doubling versus origination.
\end{definition}

\section{The ECL formula}

For a loan with amortised-cost balance $E_t$ at the start of period $t$,
marginal one-year default probability $m_t$ (conditional on survival to
$t$), and loss given default $\LGD$:

\begin{equation}
  \boxed{\;
  \ECL=\sum_{t=1}^{T} m_t\times\LGD\times E_t\times(1+\EIR)^{-t}}
  \label{eq:ecl}
\end{equation}

Three properties are prescribed by IFRS~9.5.5.17 and deserve emphasis:

\begin{enumerate}
\item \textbf{Probability-weighted and forward-looking.} The ECL must
      be a probability-weighted average over multiple macroeconomic
      scenarios --- not a single ``most likely'' path, and not a worst
      case.
\item \textbf{Discounted at the original EIR.} The allowance is a
      present value; a loss expected in year~3 is discounted at the
      loan's own effective interest rate.
\item \textbf{Marginal, not cumulative, PDs.} The $t$-th term uses the
      probability of defaulting \emph{in} period $t$ given survival
      until then. With survival factors
      $s_t=\prod_{j\le t}(1-m_j)$, the year-$t$ term can equivalently
      be written with cumulative PDs, $(S_t-S_{t-1})$, where
      $S_t=1-s_t$.
\end{enumerate}

\begin{remark}
The triad $m_t\times\LGD\times E_t$ is the same credit-risk arithmetic
as the regulatory $\EL=\PD\times\LGD\times\EAD$. The frameworks differ
in \emph{timing, horizon, aggregation, and discounting} --- not in the
loss building blocks. This is deliberate in this book: the same
parameters feed both models, which is exactly what the companion
dashboards demonstrate.
\end{remark}

\section{A short proof: lifetime dominates 12-month ECL}

\begin{theorem}[Monotonicity of the allowance in the horizon]
\label{thm:lt}
For a loan with remaining life $T\ge 2$ and non-negative marginal PDs,
\begin{equation}
  \text{lifetime ECL} \;\ge\; \text{12-month ECL},
\end{equation}
with equality iff $m_t=0$ for all $t\ge 2$ (or $T\le 1$).
\end{theorem}

\begin{proof}
The 12-month ECL is the $t=1$ term of the sum \eqref{eq:ecl}; the
lifetime ECL is the sum of the terms $t=1,\dots,T$. Every term is a
product of non-negative factors, so adding terms $t\ge 2$ cannot
decrease the total; it strictly increases it if any $m_t>0$ for
$t\ge 2$.
\end{proof}

\begin{corollary}
Stage migration from stage~1 to stage~2 never decreases the allowance
of a multi-period performing loan. This is the analytical origin of
the P\&L cliff effect.
\end{corollary}

\section{Worked example: the \EUR{10{,}000} loan}

Loan terms: \EUR{10{,}000}, 3-year annuity, EIR $8.5\%$ (schedule in
Table~\ref{tab:amortisation}). Marginal PDs (base scenario,
conditional on survival): $m=(2.0\%,\,2.5\%,\,3.0\%)$;
$\LGD=45\%$.

\subsection{Stage 1: 12-month ECL}

\begin{equation}
  \text{loss year 1 (undiscounted)}=0.02\times0.45\times10{,}000
  =\EUR{90.00},
\end{equation}
\begin{equation}
  \ECL_{12\text{m}}=\frac{90.00}{1.085}=\EUR{82.95}.
\end{equation}

\subsection{Lifetime ECL, base scenario}

\begin{table}[htbp]
\centering
\caption{Lifetime ECL decomposition, base scenario.}
\label{tab:ltbase}
\begin{tabular}{rrrrr}
\toprule
Year $t$ & $m_t$ & $E_t$ & $m_t\LGD E_t$ & PV at $8.5\%$\\
\midrule
1 & 2.0\% & 10{,}000.00 & 90.00 & 82.95\\
2 & 2.5\% & 6{,}934.61 & 78.01 & 66.27\\
3 & 3.0\% & 3{,}608.66 & 48.72 & 38.14\\
\midrule
\multicolumn{4}{r}{Lifetime ECL (base)} & 187.36\\
\bottomrule
\end{tabular}
\end{table}

\subsection{Scenario weighting}

IFRS~9.5.5.17 requires probability weighting across scenarios. We use
three: base ($60\%$), optimistic with halved PDs ($20\%$), adverse with
doubled PDs ($20\%$). Because every term of \eqref{eq:ecl} is linear in
the $m_t$, each scenario's ECL scales linearly in its PD scale:
\begin{align}
  \ECL_{\text{opt}}&=93.68,\qquad \ECL_{\text{base}}=187.36,\qquad
  \ECL_{\text{adv}}=374.72,\\
  \ECL_{\text{lt}}
  &=0.2\times93.68+0.6\times187.36+0.2\times374.72
   =\EUR{206.10}.
\end{align}
The same weighting applied to the stage-1 horizon gives the
probability-weighted 12-month ECL
$0.2\times 41.47+0.6\times82.95+0.2\times165.90=\EUR{91.24}$
(the companion dashboard displays these weighted figures).

\begin{table}[htbp]
\centering
\caption{Scenario-weighted lifetime ECL under alternative weightings.}
\label{tab:weights}
\begin{tabular}{lrrr}
\toprule
Weighting (opt/base/adv) & Optimistic share & Adverse share & ECL\\
\midrule
20/60/20 (baseline) & 20\% & 20\% & 206.10\\
33/34/33 & 33\% & 33\% & 218.27\\
10/80/10 & 10\% & 10\% & 196.73\\
50/0/50 & 50\% & 50\% & 234.20\\
\bottomrule
\end{tabular}
\end{table}

\keyinsight{The weighting is linear in the scenario PDs, so the
\emph{choice of weights} matters less than the \emph{shape of the
scenario fan}. The non-linearity of real ECL models comes from PD
responses to macro drivers inside each scenario, not from the averaging
step.}

\subsection{Stage-2 migration: the cliff in numbers}

Suppose that after one year the borrower's credit risk has
deteriorated such that the remaining lifetime PD has roughly doubled
(a plausible SICR trigger). The loan migrates from stage~1 to
stage~2. For the remaining periods the allowance becomes the lifetime
figure on the then-outstanding balance of \EUR{6{,}934.61} --- about
\EUR{104} for years 2--3 ($66.27+38.14$, base scenario) versus
\EUR{82.95} under a fresh 12-month view. Nothing about the contractual
cash flows has changed; only the \emph{expected timing} of losses. A
bank with a large book migrating simultaneously books a large
impairment charge --- the mechanism behind the P\&L shocks of 2020.

\subsection{The discounting subtlety}

Because the ECL is discounted at the loan's own EIR, a \emph{higher}
contract rate lowers the discounted allowance:
\begin{center}
\begin{tabular}{lcccc}
\toprule
EIR & 5\% & 8.5\% & 12\% & 20\%\\
\midrule
12-month ECL & 85.71 & 82.95 & 80.36 & 75.00\\
\bottomrule
\end{tabular}
\end{center}
All else equal, a riskier-priced loan (higher EIR) shows a slightly
\emph{smaller} allowance for the same expected loss --- a real,
disclosed, and sometimes counter-intuitive effect of the standard.

\section{Reconciliation with Part A}

\begin{itemize}
\item Both frameworks use PD, LGD, and an exposure notion. The
      regulatory year-1 expected loss (\EUR{90}) equals the
      undiscounted first-period IFRS~9 loss in the base scenario ---
      deliberately, because the same parameters feed both models.
\item Regulation holds capital against the \emph{tail}
      ($K\times\EAD=\EUR{463.89}$ covering the 99.9\%\ conditional
      loss \emph{beyond} expectation); IFRS~9 books the
      \emph{expectation} (\EUR{82.95}--\EUR{206.10}) as a provision.
      Both coexist: an A-IRB bank books IFRS~9 ECL as a provision and
      additionally holds IRB capital.
\item Under CRR Arts.~158 and 160--162, the excess of IRB expected
      losses over accounting provisions (or vice versa) enters the
      own funds calculation. The details are bank-specific and outside
      our scope.
\end{itemize}

\section{Limitations of Part B}

\begin{itemize}
\item Real banks run score-to-PD models with macroeconomic overlays and
      dozens of scenarios; three scenarios are the minimum acceptable
      teaching device.
\item SICR assessment is partly judgemental; the ``lifetime PD doubling''
      rule is market practice, not a standard requirement.
\item IFRS~9 exposure is the amortised-cost carrying amount (plus, for
      revolvers, expected drawdowns), which need not equal the
      regulatory EAD.
\item Stage~3 (credit-impaired) measurement --- the probability-weighted
      shortfall of cash flows --- is not modelled here.
\item A further subtlety: marginal PDs conditioned on survival make the
      ECL depend on the \emph{whole} PD term structure, not just on the
      first-year PD.
\end{itemize}

% =====================================================================
\chapter{Part C --- Risk-Sensitive Pricing}
\label{chap:pricing}
% =====================================================================

\section{The pricing question}

Given the risk of the loan (Part~A) and its expected loss and
provisioning (Part~B), what interest rate should the bank charge so
that the loan earns its cost of capital? This is risk-based pricing:
prices that reflect the borrower's risk profile --- standard practice
in consumer and mortgage lending.

\section{The RAROC framework}

Risk-Adjusted Return on Capital (RAROC), developed at Bankers Trust in
the late 1970s, allocates the loan's risk capital --- here the
regulatory own funds requirement --- and compares the loan's
risk-adjusted profit with the shareholders' hurdle rate (cost of
equity, COE):

\begin{equation}
  \raroc=\frac{\text{revenue}-\text{cost of funds}
  -\text{operating cost}-\text{expected loss}}{\text{capital}} .
  \label{eq:rarocdef}
\end{equation}

\begin{definition}[Pricing rule]
Accept or price the loan such that $\raroc\ge \text{COE}$. The
\emph{break-even rate} is the customer rate $r_L$ that makes
$\raroc=\text{COE}$ exactly.
\end{definition}

\section{Derivation of the required rate}

Let $B$ be the average outstanding balance over the life of the loan,
$r_L$ the customer rate, $r_F$ the funding rate, $O$ the amortised
operating cost per year, $\EL_{\text{p.a.}}$ the expected loss charged
to pricing per year, $C=0.08\times\RWA$ the allocated capital, and $h$
the hurdle rate. Setting \eqref{eq:rarocdef} equal to $h$:
\begin{align}
  (r_L-r_F)\,B - O - \EL_{\text{p.a.}} &= h\,C\\
  \Aboxed{r_L &= r_F+\frac{O+\EL_{\text{p.a.}}+h\,C}{B}}
  \label{eq:breakeven}
\end{align}

Equation~\eqref{eq:breakeven} is the \emph{additive pricing formula}:
customer rate $=$ funding cost $+$ operating cost $+$ expected loss
$+$ cost of capital, each per unit of exposure. The third term
converts Part~B into price; the fourth converts Part~A into price.

\section{Worked example: pricing the \EUR{10{,}000} loan}

Assumptions (illustrative, per year unless stated):
\begin{itemize}
\item funding rate $r_F=2.0\%$;
\item one-time origination/servicing cost \EUR{500}, amortised over
      3 years: $O=\EUR{166.67}$ p.a.;
\item expected loss on the 12-month basis:
      $\EL_{\text{p.a.}}=\PD\,\LGD\,\EAD=\EUR{90.00}$;
\item capital $C=\EUR{463.89}$ (Part~A); hurdle rate $h=12\%$;
\item average outstanding balance over 3 years (from
      Table~\ref{tab:amortisation}):
      $B=(10{,}000+6{,}934.61+3{,}608.66)/3=\EUR{6{,}847.75}$.
\end{itemize}

\begin{equation}
  r_L = 0.02+\frac{166.67+90.00+0.12\times 463.89}{6{,}847.75}
      = 0.02+\frac{312.34}{6{,}847.75}
      = 6.56\% \;\text{p.a.}
\end{equation}

\begin{table}[htbp]
\centering
\caption{Decomposition of the $4.56\%$ risk-and-cost margin.}
\label{tab:decomp}
\begin{tabular}{lrl}
\toprule
Component & Margin share & Basis\\
\midrule
Operating cost & 2.43\% & $166.67/6{,}847.75$\\
Expected loss & 1.31\% & $90.00/6{,}847.75$\\
Cost of capital & 0.81\% & $0.12\times463.89/6{,}847.75$\\
\midrule
Total margin & 4.56\% & \\
\bottomrule
\end{tabular}
\end{table}

\section{Cross-check: RAROC at the example EIR of 8.5\%}

If the bank charges $8.5\%$ (the rate assumed for the amortisation
schedule in Part~B), the RAROC is
\begin{equation}
  \raroc(8.5\%)
  =\frac{(0.085-0.02)\times 6{,}847.75-166.67-90.00}{463.89}
  =\frac{188.44}{463.89}=40.6\% .
\end{equation}

A $40\%$ RAROC looks implausibly high --- and this is the model's most
important lesson: \emph{regulatory capital for a low-PD retail loan is
small}, so a RAROC computed on regulatory capital is easily inflated.
Professional pricing desks therefore use economic capital (a
bank-specific tail measure, often several times the regulatory figure
for undiversified retail books), Pillar~2 add-ons, MREL and
leverage-ratio constraints, and acquisition costs that dwarf
servicing costs. The break-even logic \eqref{eq:breakeven} is unchanged;
only the size of $C$ and $O$ grows. Table~\ref{tab:capgrid} shows the
effect of the two most important corrections.

\begin{table}[htbp]
\centering
\caption{Break-even rate (\%) as a function of hurdle rate and capital
multiplier ($\PD=2\%$, 12-month and lifetime EL basis).}
\label{tab:capgrid}
\begin{tabular}{lcccccc}
\toprule
& \multicolumn{3}{c}{12-month EL basis} & \multicolumn{3}{c}{lifetime EL basis}\\
\cmidrule(lr){2-4}\cmidrule(lr){5-7}
Capital multiplier & $h=8\%$ & $h=12\%$ & $h=16\%$
                  & $h=8\%$ & $h=12\%$ & $h=16\%$\\
\midrule
1.0 (regulatory) & 6.29 & 6.56 & 6.83 & 5.98 & 6.25 & 6.52\\
2.0              & 6.83 & 7.37 & 7.92 & 6.52 & 7.06 & 7.60\\
4.0 (econ.\ approx.)& 7.92 & 9.00 & 10.08 & 7.60 & 8.69 & 9.77\\
\bottomrule
\end{tabular}
\end{table}

\keyinsight{With economic capital of four times regulatory capital and
a $12\%$ cost of equity, the teaching model's break-even rate
($9.00\%$) lands in the region of observed market rates for unsecured
consumer loans --- the model closes the gap between regulatory
arithmetic and observed prices. Use the multiplier slider in
\texttt{Pricing\_visualised.py} to explore this.}

\section{Risk sensitivity: two borrowers, same product}

Take a riskier borrower with $\PD=8\%$, $\LGD=45\%$ (row of
Table~\ref{tab:pdsens}): $R=3.79\%$, $\PD^{\text{stress}}=20.64\%$,
$K=5.69\%$, $RW=71.1\%$, $\RWA=\EUR{7{,}108}$, own funds
\EUR{568.64}, and an EL of $0.08\times0.45=\EUR{360}$ p.a. The
required rate becomes
\begin{equation}
  r_L=0.02+\frac{166.67+360.00+0.12\times 568.64}{6{,}847.75}
  =10.69\%\;\text{p.a.},
\end{equation}
against $6.56\%$ for the $2\%$ borrower. Note the composition:
regulatory capital rises only moderately (the falling correlation
cushions the stress PD), while the \emph{expected loss} term does most
of the pricing work. Risk-based pricing in retail is, in this precise
sense, mostly expected-loss pricing plus a modest capital charge.

\section{Pricing with a lifetime-EL view}

Part~B showed the probability-weighted lifetime ECL
(\EUR{206.10}) is a one-off present value, while the recurring
12-month EL charged every year (\EUR{90} p.a.\ $\times$ 3 years,
ignoring run-off) overstates total expected losses. If management
prices off the amortised lifetime view, the EL term becomes
$206.10/3=\EUR{68.70}$ p.a., lowering the required rate to about
$6.25\%$ (Table~\ref{tab:capgrid}, right panel). The choice of EL
basis --- recurring 12-month versus amortised lifetime --- is a
genuine, disclosed pricing-policy decision, not a technicality: the
12-month basis is the conservative one because it charges a full-EAD
expected loss even as the balance amortises.

\section{Limitations of Part C}

\begin{itemize}
\item RAROC on regulatory capital is biased upward for low-risk
      retail; economic capital, RARORAC, or EVA are the professional
      fixes.
\item The model ignores taxes, the liquidity cost of the funding base,
      prepayment risk (consumer loans prepay), cross-selling
      subsidies, adverse-selection feedback (raising the price worsens
      the retained pool), and regulatory price constraints.
\item A single hurdle rate ignores portfolio context: a new loan is
      priced on \emph{marginal} capital, not average capital.
\end{itemize}

% =====================================================================
\chapter{Synthesis: One Loan, Three Numbers}
% =====================================================================

\begin{table}[htbp]
\centering
\caption{The running example across the three frameworks (base case).}
\begin{tabular}{lll}
\toprule
Question & Framework & Number\\
\midrule
Unexpected tail-loss capital & CRR A-IRB (Part A)
  & \EUR{463.89} capital ($\RWA=\EUR{5{,}798.64}$)\\
Expected-loss provision & IFRS 9 (Part B)
  & \EUR{82.95} (stage 1) to \EUR{206.10} (stage 2, weighted lifetime)\\
Break-even customer rate & RAROC pricing (Part C)
  & $6.56\%$ p.a.\ (12-month EL basis)\\
\bottomrule
\end{tabular}
\end{table}

The three numbers are mechanically linked: the same $\PD$, $\LGD$, and
exposure drive all three. Change the PD from $2\%$ to $8\%$ and all
three move --- the capital requirement rises by about a quarter
(correlation cushioning), the stage-1 allowance roughly quadruples,
and the required rate rises by about two thirds, driven mainly by
expected loss. This coupling is exactly what the three companion
dashboards visualise interactively.

\begin{figure}[htbp]
\centering
\begin{tikzpicture}[node distance=8mm and 12mm,
  box/.style={draw=ceruleanlink, rounded corners, fill=white,
              minimum width=30mm, align=center}]
\node[box] (pd) {$\PD=2\%$\\$\LGD=45\%$\\$\EAD=10{,}000$};
\node[box, right=of pd] (a) {Part A\\capital \EUR{463.89}\\(99.9\%\ tail)};
\node[box, right=of a] (b) {Part B\\allowance\\$\EUR{82.95}$--$\EUR{206.10}$};
\node[box, below=8mm of b] (c) {Part C\\break-even $6.56\%$};
\draw[-{Latex[length=2mm]}] (pd) -- (a);
\draw[-{Latex[length=2mm]}] (pd) -- (b);
\draw[-{Latex[length=2mm]}] (a.south) |- (c.west)
  node[near start, right, font=\small] {$C$};
\draw[-{Latex[length=2mm]}] (b) -- (c) node[midway, right, font=\small] {$\EL$};
\end{tikzpicture}
\caption{One parameter set drives capital, allowance, and price.}
\end{figure}

% =====================================================================
\chapter{The Companion Dashboards}
\label{chap:dashboards}
% =====================================================================

The three Python scripts implement the mathematics of this book as
self-contained Flask dashboards (Python 3.12; dependencies per
\texttt{requirements\_py312.txt}; runnable in Spyder via F5).

\begin{center}
\begin{tabular}{llp{78mm}}
\toprule
Script & Port & Demonstrates\\
\midrule
\texttt{RWA\_visualised.py} & 5001 &
risk weight vs PD; $K$ and $R$ vs PD; RWA vs EAD;
reproduces Tables~\ref{tab:pdsens} and \ref{tab:lgdsens}.\\
\texttt{Impairment\_visualised.py} & 5002 &
stage-1 vs stage-2 allowance; scenario weighting; per-year ECL
decomposition; reproduces Tables~\ref{tab:ltbase} and \ref{tab:weights}.\\
\texttt{Pricing\_visualised.py} & 5003 &
break-even rate decomposition; RAROC vs candidate rate with hurdle
line; capital multiplier; EL basis switch; reproduces
Table~\ref{tab:capgrid}.\\
\bottomrule
\end{tabular}
\end{center}

Recommended study sequence: reproduce Section~\ref{sec:workedA} in
\texttt{RWA\_visualised.py}, then vary the PD slider and compare with
Table~\ref{tab:pdsens}; move to
\texttt{Impairment\_visualised.py} and switch the stage selector to
see the cliff; finish with \texttt{Pricing\_visualised.py}, first with
the capital multiplier at 1.0, then at 4.0.

% =====================================================================
\chapter{Exercises}
% =====================================================================

\begin{enumerate}
\item Derive the factor $12.5$ from the $8\%$\ rule without looking at
      Chapter~\ref{chap:rwa}. Explain in one sentence why the factor
      contains no economics.
\item Compute $R(\PD)$ for $\PD=4\%$ by hand (two decimal places) and
      verify with \texttt{RWA\_visualised.py}.
\item Show algebraically that the two expressions
      $\Phi\!\bigl[(G(\PD)+\sqrt R\,G(0.999))/\sqrt{1-R}\bigr]$ and
      \eqref{eq:pdstress} are identical.
\item A supervisor raises the stress confidence level from
      $99.9\%$ to $99.99\%$. Explain qualitatively (or compute) what
      happens to $K$ for our running example.
\item Compute the stage-1 and stage-2 allowances of the running
      example with $m=(3\%,3.5\%,4\%)$ and interpret the change versus
      Table~\ref{tab:ltbase}.
\item Prove or disprove: shifting weight from the base to the adverse
      scenario always increases the weighted ECL.
\item Using Table~\ref{tab:capgrid}, estimate the break-even rate for
      capital multiplier 3.0 and $h=14\%$ (interpolate), then verify
      with \texttt{Pricing\_visualised.py}.
\item Explain why the RAROC of the running example at a customer rate
      of $8.5\%$ exceeds $40\%$, and why this does not imply that
      retail lending earns a $40\%$ return on equity in practice.
\end{enumerate}

\paragraph{Hints.} (1) Pure algebra: division by $0.08$. (3) Split the
numerator over $\sqrt{1-R}$. (4) $G(0.9999)\approx 3.719$ instead of
$3.090$ --- the stress PD rises further. (6) True --- by linearity in
the adverse scenario's weight, since the adverse ECL exceeds the base
ECL.

% =====================================================================
\appendix
\chapter{Notation}
\label{app:notation}
% =====================================================================

\begin{center}
\begin{tabular}{ll}
\toprule
Symbol & Meaning\\
\midrule
$\PD$ & probability of default (one year)\\
$\LGD$ & loss given default\\
$\EAD$ & exposure at default\\
$R$ & asset correlation (systematic weight squared)\\
$Y$, $\varepsilon$ & systematic / idiosyncratic standard normal factors\\
$\Phi$, $G=\Phi^{-1}$ & standard normal CDF / quantile function\\
$K$ & regulatory capital factor per unit of EAD\\
$RW$ & risk weight, $12.5\times K$\\
$\RWA$ & risk-weighted assets\\
$T$ & remaining life (years)\\
$m_t$ & marginal (conditional-on-survival) one-year PD in year $t$\\
$E_t$ & amortised-cost balance at the start of year $t$\\
$\EIR$ & effective interest rate\\
$\ECL$ & expected credit loss (allowance)\\
$\EL$ & expected loss\\
$B$ & average outstanding balance over the loan life\\
$r_L$, $r_F$ & customer rate / funding rate\\
$C$ & allocated (regulatory or economic) capital\\
$h$ & hurdle rate (cost of equity)\\
$\raroc$ & risk-adjusted return on capital\\
\bottomrule
\end{tabular}
\end{center}

% =====================================================================
\chapter{Source Notes and References}
% =====================================================================

\section*{References}

\begin{enumerate}[label={[\arabic*]}]
\item Regulation (EU) No 575/2013 (CRR), Art.~153 (risk-weight
      functions, corporates) and Art.~154 (retail exposures,
      correlation formula). EUR-Lex:
      \url{https://eur-lex.europa.eu/eli/reg/2013/575/oj/eng};
      Art.~154 text:
      \url{https://www.legislation.gov.uk/eur/2013/575/article/154}.
\item Basel Committee on Banking Supervision (2005), ``An Explanatory
      Note on the Basel II IRB Risk Weight Functions'', BIS:
      \url{https://www.bis.org/bcbs/irbriskweight.pdf}.
\item Gordy, M.~B. (2003), ``A risk-factor model foundation for
      ratings-based bank capital rules'', \emph{Journal of Financial
      Intermediation} 12, 199--232.
\item Vasicek, O. (2002), ``The Distribution of Loan Portfolio
      Value'', \emph{Risk} 15(12), 160--162; Merton, R.~C. (1974),
      ``On the Pricing of Corporate Debt'', \emph{Journal of Finance}
      29, 449--470.
\item BIS Financial Stability Institute, ``IFRS~9 and expected loss
      provisioning --- Executive Summary'':
      \url{https://www.bis.org/fsi/fsisummaries/ifrs9.pdf};
      IFRS~9 \emph{Financial Instruments}, IFRS Foundation
      (ifrs.org), paras 5.5.15--5.5.20.
\item Risk-adjusted return on capital (RAROC):
      \url{https://en.wikipedia.org/wiki/Risk-adjusted_return_on_capital};
      risk-based pricing:
      \url{https://en.wikipedia.org/wiki/Risk-based_pricing}.
\item User-provided materials: \texttt{old\_chat.txt} (prior worked
      derivation of the RWA example) and
      \texttt{requirements\_py312.txt} (target Python environment).
\end{enumerate}

All web sources accessed October 2026. Credibility ratings follow in
the Source Notes table.

\section*{Source Notes}

\begin{center}
\begin{tabular}{llc}
\toprule
Source & Credibility & Last updated\\
\midrule
EUR-Lex, Regulation (EU) 575/2013 & 5/5 & 2013 (consolidated ongoing)\\
Legislation.gov.uk, CRR Art.~154 & 5/5 & 2013 (EU-exit amendments noted)\\
BIS, Explanatory Note on Basel II IRB Functions & 5/5 & 2005\\
BIS FSI, IFRS 9 Executive Summary & 5/5 & 2017\\
Wikipedia, RAROC & 3/5 & ---\\
Wikipedia, Risk-based pricing & 3/5 & ---\\
\bottomrule
\end{tabular}
\end{center}

\textbf{Caveats.} Article numbering and the correlation constants were
verified against the retrieved CRR Art.~154 text; the CRR2 output-floor
timeline and CRR3 changes were not re-verified against the current
consolidated text and are stated qualitatively only. The IFRS~9
paragraph references follow the IFRS Foundation numbering and the BIS
FSI summary.

\section*{Open Questions}

\begin{enumerate}
\item What is the bank's actual (supervisor-approved) PD and LGD for
      this borrower segment? Unobservable from outside; only the bank
      can answer.
\item Economic capital for a concentrated retail book exceeds the
      \EUR{463.89} regulatory figure --- by how much requires the
      bank's internal tail model.
\item The current CRR3 output-floor percentage applicable to the
      bank's IRB portfolio changes effective RWA; not modelled here.
\item The servicer's actual stage-2 migration rules (days-past-due
      triggers, cure rules) materially affect IFRS~9 P\&L timing.
\end{enumerate}

\section*{Recommendations and Next Steps}

\begin{enumerate}
\item Run the three companion scripts and reproduce every table in
      this book; each script prints its parameter set and results.
\item Stress the parameters: PD up to $20\%$, LGD to $90\%$, and
      observe the concavity of the risk weight and the near-linearity
      of the ECL.
\item Read the primary texts: CRR Arts.~153--154 and the BIS
      explanatory note for Part~A; IFRS~9 paras 5.5.15--5.5.20 for
      Part~B.
\item Extend the pricing script to economic capital ($3$--$5\times$
      regulatory capital for retail) and observe the break-even rate
      converge toward observed market rates.
\end{enumerate}

\vfill
\begin{center}
\color{accentgrey}\small
This textbook was created with AI assistance; all numerical examples
were computed and cross-checked with executable code. It is an
educational treatment with didactic simplifications, not regulatory,
accounting, or financial advice.
\end{center}

\end{document}
